\section{Predictable boundary inputs preserve the uniform moment exponents}
\label{cm:section}

Let $\mathcal Q$ be the class of nonanticipating input policies which add
a leaf at the retained root at predictable intensity $q_t\in[0,1/2]$.
Apart from these additions, the component follows the inert generator.
The root starts in the newborn state $x_\infty$, with two real vertices.
The policies may use the observed component history and additional
randomness. An actual exterior supplies a policy in this class.

Theorem~\ref{sc:window} implies that,
for every $\alpha>\lambda$, there is a function $U$ and finite positive
constants $m,A,J$ such that
\begin{equation}\label{cm:window-input}
 mC\le U\le AC,\qquad
 (\alpha-\mathcal L_\infty)U\ge C,\qquad
 |U(y')-U(y)|\le J
\end{equation}
for every ordinary leaf addition $y\to y'$. In particular, one may use
$U_{\alpha,T}$ from that construction, with
$m=(1-e^{-\alpha T})/[\alpha(1-q)]$ and its already stated $A,J$.
Here the scalar $q<1$ is the finite-window contraction parameter.

\begin{theorem}[All controlled moments have the same normalized rate]
\label{cm:all-moments}
For every real $p\ge1$ and every $a>\lambda$, there is a finite constant
$K_{p,a}$, independent of the input policy, such that
\begin{equation}\label{cm:bound}
 \sup_{q\in\mathcal Q}\mathbb E_q C(t)^p
 \le K_{p,a}e^{pa t}\qquad(t\ge0).
\end{equation}
Consequently,
\begin{equation}\label{cm:exponent}
 \boxed{\lim_{t\to\infty}\frac1t
 \log\sup_{q\in\mathcal Q}\mathbb E_q C(t)^p=p\lambda.}
\end{equation}
For computable parameters and certified positive slack, the constants
are effective once a finite-window certificate at any
$\alpha\in(\lambda,a)$ has been supplied. More generally, the same
moment upper bound holds from any fixed finite initial component, with
constants depending on that component.
\end{theorem}
\begin{proof}
Choose $\lambda<\alpha<a$. The established lower bound on $\mu$ gives
$\alpha>\lambda\ge\mu\ge0.565>1/2$. Set
\[
 W=U+\frac{J}{2\alpha}.
\]
For every admissible policy, its conditional generator satisfies
\begin{equation}\label{cm:linear-drift}
 (\mathcal L_\infty+q_t\Delta_{\rm root})W
 \le\alpha U-C+J/2\le\alpha W.
\end{equation}
The jumps of $W$ have absolute size at most $J$, and $W\ge mC\ge m$.

Let $N_t$ be the real vertex count and put $Y_t=N_t+1/2$. Its birth
intensity is at most $Y_t$. For each integer $r\ge1$, the rising
factorial $y^{\overline r}=y(y+1)\cdots(y+r-1)$ satisfies
\[
 y\bigl[(y+1)^{\overline r}-y^{\overline r}\bigr]
 =r y^{\overline r}.
\]
It follows by localization that
\begin{equation}\label{cm:population}
 \mathbb E Y_t^r\le
 \mathbb E Y_t^{\overline r}
 \le Y_0^{\overline r}e^{rt}.
\end{equation}

First take an even integer $p=2r\ge2$. Taylor's formula and
$|\Delta W|\le J$ give, at every transition,
\[
 (W+\Delta W)^p-W^p-pW^{p-1}\Delta W
 \le\frac{p(p-1)}2J^2(W+J)^{p-2}.
\]
Since $W\ge m$, set
$K_p=\tfrac12p(p-1)J^2(1+J/m)^{p-2}$. Summing over transitions and
using \eqref{cm:linear-drift} yields
\[
 (\mathcal L_\infty+q_t\Delta_{\rm root})W^p
 \le p\alpha W^p+K_p W^{p-2}Y.
\]
For $p>2$, Young's inequality, or maximization in $W$, gives for every
$\eta>0$ a finite explicit $B_{p,\eta}$ with
\[
 K_p W^{p-2}Y\le\eta W^p+B_{p,\eta}Y^{p/2}.
\]
For example, one may take
\[
 B_{p,\eta}=\frac2p K_p
 \left(\frac{(p-2)K_p}{p\eta}\right)^{(p-2)/2}.
\]
For $p=2$, the corresponding remainder is simply $J^2Y$.
Choose $\eta=p(a-\alpha)$ when $p>2$; in the $p=2$ case increase
the drift coefficient $2\alpha$ harmlessly to $2a$. Thus
\[
 \frac d{dt}\mathbb E W_t^p
 \le pa\mathbb E W_t^p+B\,Y_0^{\overline {p/2}}e^{pt/2},
\]
with $B=B_{p,\eta}$ for $p>2$ and $B=J^2$ for $p=2$.
Because $a>1/2$, integration gives
\begin{equation}\label{cm:explicit-even}
 \mathbb E W_t^p
 \le e^{pa t}\left[
 W_0^p+\frac{B\,Y_0^{\overline {p/2}}}{p(a-1/2)}\right].
\end{equation}
Divide by $m^p$. All localization limits are justified by
$W\le AC+J/(2\alpha)\le AN+J/(2\alpha)$ and the finite-horizon
population moments. These estimates hold uniformly for every policy.

For a general real $p\ge1$, choose an even integer $k\ge p$ and use
$\mathbb E C^p\le(\mathbb E C^k)^{p/k}$. This proves
\eqref{cm:bound}. Conversely the controlled class contains every admissible
newborn process, including inert exterior profiles. Jensen's inequality gives
\[
 \sup_{q\in\mathcal Q}\mathbb E_q C(t)^p
 \ge \left(\sup_x\mathbb E_x C_{\rm tag}(t)\right)^p=G(t)^p.
\]
The established limit $t^{-1}\log G(t)\to\lambda$, followed by
$a\downarrow\lambda$ in the upper bound, proves
\eqref{cm:exponent}.
\end{proof}

